\documentclass{article}
\usepackage{amsmath, amssymb, amsthm}
\usepackage{hyperref}
\usepackage{geometry}
\geometry{margin=1in}

\title{Erd\H{o}s Problem \#950}
\date{Last updated: 2025-08-31}
\author{Source: erdosproblems.com}

\begin{document}
\maketitle

\noindent\textbf{Status:} OPEN \quad \textbf{No prize}

\noindent\textbf{Tags:} number theory, primes

\medskip

\noindent\textbf{Problem Statement:}

Let\[f(n) = \sum_{p&#60;n}\frac{1}{n-p}.\]Is it true that\[\liminf f(n)=1\]and\[\limsup f(n)=\infty?\]Is it true that $f(n)=o(\log\log n)$ for all $n$?




This function was considered by de Bruijn, Erd\H{o}s, and Tur\'{a}n, who showed that\[\sum_{n&lt;x}f(n)\sim \sum_{n&lt;x}f(n)^2\sim x.\]They give no proofs, but a proof of the (harder) second claim is given by Gorodetsky here . The existence of some $c&gt;0$ such that there are $\gg n^c/\log n$ many primes in $[n,n+n^c]$ implies that $\liminf f(n)&gt;0$. Erd\H{o}s writes that a 'weaker conjecture which is perhaps not quite inaccessible' is that, for every $\epsilon&gt;0$, if $x$ is sufficiently large there exists $y&lt;x$ such that\[\pi(x)&lt; \pi(y)+\epsilon \pi(x-y).\](Compare this to [855] .) He notes that if\[\pi(x)&lt; \pi(y)+O\left(\frac{x-y}{\log x}\right)\]for all $y&lt;x-(\log x)^C$ for some constant $C&gt;0$ then $f(n)\ll \log\log\log n$. The study of $f(p)$ is even harder, and Erd\H{o}s could not prove that\[\sum_{p&lt;x}f(p)^2\sim \pi(x).\]See also [1210] .


\bigskip
\noindent\small{Source: \url{https://www.erdosproblems.com/950}}
\end{document}
